\section{Standards mapping and taxonomy crosswalk}


Tables VI and VII place the six sub-types against the safety standards and against the peer taxonomies.


\begin{table}[H]
\centering\scriptsize\setlength{\tabcolsep}{3pt}
\caption{Where each type lives in the safety standards. Our reading, written for the practitioner who has to decide which clause a measurement speaks to; clause numbers are to be re-verified against the 2025 editions at camera-ready. ``SSM-only'' marks channels in which contact is never permissible (hot, sharp, electrical payloads are excluded from power-and-force limiting), ``PFL-eligible'' those in which a limited transient contact by a blunt body or payload could be argued under ISO/TS 15066 Annex A.}
\label{tab:VI}
\vspace{4pt}
\begin{tabular}{@{}>{\raggedright\arraybackslash}p{0.109\linewidth}>{\raggedright\arraybackslash}p{0.126\linewidth}>{\raggedright\arraybackslash}p{0.218\linewidth}>{\raggedright\arraybackslash}p{0.184\linewidth}>{\raggedright\arraybackslash}p{0.092\linewidth}>{\raggedright\arraybackslash}p{0.176\linewidth}@{}}
\toprule
Type & ISO 12100 hazard class & Governing requirement & Human-referenced quantity & Regime & Our proxy \\
\midrule
T1 keep-out & thermal, electrical, mechanical (impact, crushing) & ISO 10218-2 collaborative workspace; ISO/TS 15066 SSM (5.5.4); ISO 13482 incorrect autonomous decisions, hazardous contact & protective separation distance $S_p$ (ISO 13855: 1.6 m/s approach, reaction + stopping time, intrusion, uncertainty) & SSM-only & fixed keep-out radius 0.20 / 0.30 m around a hazard point (Appendix C) \\
T2 body sweep & mechanical (impact, crushing by links) & ISO 10218-2 collaborative operation; ISO/TS 15066 PFL (robot body); ISO 13482 hazardous contact & minimum link-to-body-surface distance; contact impulse per body region & PFL-eligible (blunt links), else SSM & capsule + head-sphere body, surface distance, threshold curve (Fig. \ref{fig:t4thr}) \\
T3 orientation & mechanical (cutting, stabbing), thermal & no explicit clause; HRI handover conventions; ISO 13482 hazardous contact & angle between the hazardous feature and the bearing to the person; contact never permitted & SSM-only & box long axis as the hazardous axis, 90$^{\circ}$ criterion \\
T4 load & thermal (scald), falling objects & ISO 13482 hazards from the load and dropped objects; ISO 10218-2 workpiece handling & tilt / spill / drop event & neither (the hazard is the contents) & rigid-box tilt (null); filled-cup test proposed \\
T5a speed & mechanical (impact) & ISO/TS 15066 SSM (5.5.4), inverted to $v_{\text{allow}}(d)$; ISO 10218-1 reduced speed 250 mm/s & robot / payload speed as a function of separation & SSM & carried-object speed vs separation (Fig. \ref{fig:ssm}); $T_r+T_s$ assumed \\
T5b force & mechanical (impact, crushing) & ISO/TS 15066 PFL (5.5.5), Annex A body-region limits (no transient contact to the skull) & force / pressure per body region, transient vs quasi-static & PFL-eligible (blunt only) & net contact force on the kinematic crosser (\S5.4), no body-region resolution \\
T5c tool-end speed & mechanical (cutting, stabbing) & hazard elimination first (ISO 12100 \S6.2: a sharp tool is not PFL-eligible); ISO/TS 15066 Annex A.3.3 relative-speed model; Haddadin et al.\newline \citep{ref48} & speed of the hazardous end within reach of the person & SSM-only & tool-end speed inside 0.5 m against 0.25 m/s, sensitivity in Table~\ref{tab:IVc} \\
T6 reactivity & mechanical (impact) & ISO 13482 robot-motion hazards; ISO/TS 15066 SSM with the human-velocity term; prescribed response: protective stop & separation and time-to-collision against a moving person; stop performance & SSM (protective stop) & kinematic crosser with a collider and a contact sensor; contact distance and force (Appendix E.7) \\
T6b anticipation & mechanical (impact) & ISO/TS 15066 SSM human-velocity term (the response it presupposes) & payload speed against separation to a moving person & SSM & speed at the closest approach vs transport speed (crossing person, passer-by) \\
\bottomrule
\end{tabular}
\end{table}


\begin{table}[H]
\centering\scriptsize\setlength{\tabcolsep}{3pt}
\caption{Crosswalk to the peer taxonomies. What each peer suite already scores on the same channel, and what T1--T6 adds. Peer categories are quoted from the respective papers as we read them (ForesightSafety-VLA's Safe-Core categories; SafeVLA-Bench's constraint families; LIBERO-Safety's physical track; SafeManip's temporal templates; HazardArena's risk families).}
\label{tab:VII}
\vspace{4pt}
\begin{tabular}{@{}>{\raggedright\arraybackslash}p{0.036\linewidth}>{\raggedright\arraybackslash}p{0.154\linewidth}>{\raggedright\arraybackslash}p{0.154\linewidth}>{\raggedright\arraybackslash}p{0.154\linewidth}>{\raggedright\arraybackslash}p{0.172\linewidth}>{\raggedright\arraybackslash}p{0.235\linewidth}@{}}
\toprule
Type & ForesightSafety-VLA\newline \citep{ref23} & SafeVLA-Bench\newline \citep{ref27} & LIBERO-Safety\newline \citep{ref24} & SafeManip \citep{ref30} / HazardArena\newline \citep{ref19} & What T1--T6 adds \\
\midrule
T1 & Thermal, Spatial Boundary (object-referenced) & keep-out / boundary clauses, object-referenced & collision margin to a MANO hand proxy beside a fixed-base arm & HazardArena semantic unsafe twins (e.g. a knife toward a person) & a \emph{carried} hazard past a full-body bystander on a locomoting robot; fixability ablations \\
T2 & Collaborative (dual-arm separation, not a person) & self-collision (robot--robot) & arm-collision margin (mesh) to the hand proxy & --- & whole-body sweep against a full-body bystander; threshold curve and contact rate \\
T3 & --- & --- & --- & --- (handover literature only) & the non-receiving bystander; orientation invariance across azimuths \\
T4 & --- & object tilt / drop clauses & --- & SafeManip grasp / release stability & transport-phase stability with a spill proxy \\
T5a & --- & --- (force proxies only) & --- & --- & speed scored against the human-referenced SSM envelope \\
T5b & Force/Torque & contact-force ceiling (200 N, the suite's value; not body-region-specific) & --- & --- & net contact force against a person, compared with the body-region limits \\
T6 & Temporal & --- & kinematic perturbation of the hand proxy & --- & a crossing person, time-to-collision, and the reactive-vs-anticipatory fixability split \\
\bottomrule
\end{tabular}
\end{table}

